\documentclass[11pt,letterpaper]{article}
\usepackage[margin=0.8in,headheight=14pt]{geometry}
\usepackage[T1]{fontenc}
\usepackage{lmodern,microtype,amsmath,amssymb,enumitem}
\usepackage[dvipsnames]{xcolor}
\usepackage{fancyhdr,needspace}
\usepackage[most]{tcolorbox}
\usepackage[hidelinks]{hyperref}
\definecolor{courseblue}{HTML}{1F4E79}
\definecolor{lightblue}{HTML}{EAF2F8}
\setlength{\parindent}{0pt}
\setlength{\parskip}{0.45em}
\setlist[enumerate]{leftmargin=*,topsep=3pt,itemsep=4pt,parsep=2pt}
\pagestyle{fancy}\fancyhf{}
\fancyhead[L]{\small MATH 327: Algebra 1}
\fancyhead[R]{\small Fall 2026}
\fancyfoot[C]{\thepage}
\newcommand{\problem}[1]{\par\medskip{\large\bfseries\color{courseblue}#1}\par\smallskip}
\newcommand{\partlabel}[1]{\par\textbf{(#1)}\ }
\newcommand{\Z}{\mathbb Z}
\newcommand{\R}{\mathbb R}
\newcommand{\C}{\mathbb C}
\newcommand{\GL}{\mathrm{GL}}
\newcommand{\Aut}{\operatorname{Aut}}
\newcommand{\Stab}{\operatorname{Stab}}
\newcommand{\heading}[1]{\begin{center}\colorbox{lightblue}{\parbox{0.94\textwidth}{\centering\vspace{5pt}{\Large\bfseries\color{courseblue}#1}\\[3pt]MATH 327: Algebra 1\quad |\quad Fall 2026\vspace{5pt}}}\end{center}}

\hypersetup{pdftitle={Practice Questions for Midterm 2: Solutions},pdfauthor={Manish Patnaik}}

\definecolor{boxborder}{HTML}{B8CBDC}
\definecolor{darkgray}{HTML}{333333}
\definecolor{lightyellow}{HTML}{FFF9DC}
\definecolor{yellowborder}{HTML}{E4D58C}
\newtcolorbox{solution}{
  enhanced,
  breakable,
  colback=lightyellow,
  colframe=yellowborder,
  colbacktitle=lightyellow,
  coltitle=darkgray,
  fonttitle=\bfseries,
  title={Solution},
  title after break={Solution (continued)},
  boxrule=0.6pt,
  arc=1mm,
  left=2.5mm,
  right=2.5mm,
  top=1.5mm,
  bottom=1.5mm,
  before skip=0.2em,
  after skip=1em
}


\newtcolorbox{questionblock}[1]{
 enhanced,breakable,colback=white,colframe=boxborder,
 colbacktitle=lightblue,coltitle=courseblue,fonttitle=\bfseries,
 title={Question #1},title after break={Question #1 (continued)},boxrule=0.6pt,arc=1mm,
 left=2.5mm,right=2.5mm,top=1.5mm,bottom=1.5mm,
 before skip=0.8em,after skip=0.8em
}
\begin{document}
\heading{Practice Questions for Midterm 2: Solutions}
% Full questions from the practice sheet; solutions retain their original numbering.
% Solutions 1 and 3 adapted from the instructor's Fall 2025 midterm solutions.
\Needspace{8\baselineskip}
\begin{questionblock}{1. Products and quotients}
Let $H,K\leq G$, and define $HK=\{hk:h\in H,\ k\in K\}$. Assume that $H$ normalizes $K$: $hKh^{-1}=K$ for every $h\in H$.
\begin{enumerate}[label=(\alph*)]
\item Show that $HK$ is a subgroup of $G$.
\item Show that $K\trianglelefteq HK$ and $H\cap K\trianglelefteq H$.
\item Construct an isomorphism $H/(H\cap K)\cong HK/K$. Prove that your map is well defined and is an isomorphism.
\end{enumerate}
\end{questionblock}
\begin{solution}
Assume $H,K\leq G$ and $hKh^{-1}=K$ for all $h\in H$.
\partlabel{a} The identity belongs to $HK$. For $h,h'\in H$ and $k,k'\in K$,
\[
 (hk)(h'k')=hh'\bigl(h'^{-1}kh'\bigr)k'\in HK,
\]
because $h'^{-1}kh'\in K$. Also
\[
 (hk)^{-1}=k^{-1}h^{-1}=h^{-1}\bigl(hk^{-1}h^{-1}\bigr)\in HK.
\]
Thus $HK$ is a subgroup.
\partlabel{b} If $u\in K$ and $hk\in HK$, then
\[
 (hk)u(hk)^{-1}=h(kuk^{-1})h^{-1}\in K.
\]
Applying this also to $(hk)^{-1}$ gives equality under conjugation, so $K\trianglelefteq HK$. If $u\in H\cap K$ and $h\in H$, then $huh^{-1}\in H$ because $H$ is a subgroup, and $huh^{-1}\in K$ because $H$ normalizes $K$. Applying the same argument to $h^{-1}$ proves $H\cap K\trianglelefteq H$.
\partlabel{c} Define
\[
 F:H/(H\cap K)\longrightarrow HK/K,\qquad F\bigl(h(H\cap K)\bigr)=hK.
\]
If $h(H\cap K)=h'(H\cap K)$, then $h'^{-1}h\in H\cap K\subseteq K$, so $hK=h'K$. Thus $F$ is well defined. It is a homomorphism, since
\[
 F\bigl(h(H\cap K)\,h'(H\cap K)\bigr)=hh'K=(hK)(h'K).
\]
Every coset in $HK/K$ has the form $hkK=hK$, proving surjectivity. If $F(h(H\cap K))=K$, then $h\in K$; since also $h\in H$, the original coset is $H\cap K$. Hence $F$ has trivial kernel and is injective. Therefore $F$ is the required isomorphism.
\end{solution}
\newpage
\begin{questionblock}{2. The homomorphism $S_4\to S_3$}
Let $S_4$ act on the three partitions $\Pi_1=12\mid34$, $\Pi_2=13\mid24$, and $\Pi_3=14\mid23$, where $12\mid34$ abbreviates $\{\{1,2\},\{3,4\}\}$. Let $\varphi:S_4\to S_3$ be the resulting homomorphism. You may use that $\varphi$ is surjective and that $\ker\varphi=V=\{1,(12)(34),(13)(24),(14)(23)\}$. Set $H=\varphi^{-1}(\langle(13)\rangle)$.
\begin{enumerate}[label=(\alph*)]
\item Compute $\varphi((1234))$.
\item Show that $H=\langle(1234),(12)(34)\rangle$, and determine $|H|$.
\item Identify $H$ up to isomorphism. Is $H$ normal in $S_4$? Justify both answers.
\end{enumerate}
\end{questionblock}
\begin{solution}
Use the partitions $\Pi_1=12\mid34$, $\Pi_2=13\mid24$, and $\Pi_3=14\mid23$ from the question, and set $r=(1234)$ and $s=(12)(34)$.
\partlabel{a} Applying $r$ to the entries of each partition gives
\[
 r\Pi_1=\Pi_3,\qquad r\Pi_2=\Pi_2,\qquad r\Pi_3=\Pi_1.
\]
Consequently $\boxed{\varphi(r)=(13)}$.
\partlabel{b} Write $L=\langle r,s\rangle$. Since $\varphi(r)=(13)$ and $s\in V=\ker\varphi$, both generators lie in $H$, so $L\subseteq H$. Conversely,
\[
 r^2=(13)(24),\qquad r^2s=(14)(23),
\]
so $V\subseteq L$. If $h\in H$, then $\varphi(h)$ is either $1$ or $(13)$. In the first case $h\in V\subseteq L$. In the second, $\varphi(r^{-1}h)=1$, so $r^{-1}h\in V$ and $h\in rV\subseteq L$. Thus
\[
 H=\langle(1234),(12)(34)\rangle=V\sqcup rV,\qquad \boxed{|H|=8}.
\]
\partlabel{c} Direct computation gives
\[
 r^4=s^2=1,\qquad srs^{-1}=r^{-1}.
\]
These are the dihedral relations. They define a surjective homomorphism $D_8\to H$ sending the rotation generator to $r$ and a reflection generator to $s$. Both groups have order $8$, so this map is an isomorphism: $\boxed{H\cong D_8}$.

The subgroup $H$ is \textbf{not normal} in $S_4$. Under the correspondence theorem, $H\trianglelefteq S_4$ would imply $\varphi(H)=\langle(13)\rangle\trianglelefteq S_3$. But
\[
 (123)(13)(123)^{-1}=(12)\notin\langle(13)\rangle.
\]
\end{solution}
\newpage
\begin{questionblock}{3. Conjugation in $S_3$}
Let $G$ be a group. Define $g\cdot x=gxg^{-1}$ for $g,x\in G$.
\begin{enumerate}[label=(\alph*)]
\item Show that this rule defines a left action of $G$ on itself.
\item Show that the center $Z(G)$ is the union of all orbits of cardinality $1$.
\item Show that every normal subgroup of $G$ is a union of orbits under this action.
\item Describe all orbits and stabilizers for this action when $G=S_3$.
\end{enumerate}
\end{questionblock}
\begin{solution}
\partlabel{a} For every $x\in G$, $1\cdot x=x$. Also, for $g,h,x\in G$,
\[
 (gh)\cdot x=ghx(gh)^{-1}=g(hxh^{-1})g^{-1}=g\cdot(h\cdot x).
\]
These are exactly the axioms of a left group action.
\partlabel{b} The orbit of $x$ is a singleton if and only if $gxg^{-1}=x$ for every $g\in G$. This is equivalent to $gx=xg$ for every $g$, or $x\in Z(G)$. Hence $Z(G)$ is precisely the union of the singleton orbits.
\partlabel{c} If $N\trianglelefteq G$ and $n\in N$, then $gng^{-1}\in N$ for every $g\in G$. Thus the entire orbit of each element of $N$ is contained in $N$. Since every $n$ lies in its own orbit, $N$ is a union of orbits.
\partlabel{d} The three orbits are
\[
 \{1\},\qquad \{(12),(13),(23)\},\qquad \{(123),(132)\}.
\]
The stabilizer of $x$ is its centralizer $C_{S_3}(x)$, namely
\[
\begin{array}{c|l}
 x & \Stab_{S_3}(x)\\ \hline
 1 & S_3\\
 (12)&\{1,(12)\}\\
 (13)&\{1,(13)\}\\
 (23)&\{1,(23)\}\\
 (123),\ (132)&\{1,(123),(132)\}
\end{array}
\]
For example, $1$ and $(12)$ commute with $(12)$, and its orbit has size $3$, so its stabilizer has order $6/3=2$. Likewise, a $3$-cycle has orbit size $2$, so its stabilizer has order $3$ and equals the cyclic subgroup it generates.
\end{solution}
\newpage
\begin{questionblock}{4. Conjugation in $\mathrm{GL}_2(\mathbb R)$}
Let $A=\begin{pmatrix}1&0\\0&2\end{pmatrix}$. The group $\mathrm{GL}_2(\mathbb R)$ acts on itself by conjugation: $g\cdot B=gBg^{-1}$.
\begin{enumerate}[label=(\alph*)]
\item Describe the stabilizer of $A$ under this action. Justify your answer.
\item Describe the orbit of $A$ under this action. Justify your answer.
\end{enumerate}
\end{questionblock}
\begin{solution}
Let $A=\begin{pmatrix}1&0\\0&2\end{pmatrix}$.
\partlabel{a} A matrix $P$ stabilizes $A$ if and only if $PA=AP$. For $P=\begin{pmatrix}a&b\\c&d\end{pmatrix}$,
\[
 PA=\begin{pmatrix}a&2b\\c&2d\end{pmatrix},\qquad
 AP=\begin{pmatrix}a&b\\2c&2d\end{pmatrix}.
\]
Equality forces $b=c=0$. Invertibility then requires $a,d\ne0$. Therefore
\[
 \boxed{\Stab(A)=\left\{\begin{pmatrix}a&0\\0&d\end{pmatrix}:a,d\in\R^{\times}\right\}.}
\]
\partlabel{b} Conjugation preserves the characteristic polynomial. Thus every matrix in the orbit has characteristic polynomial $(t-1)(t-2)$, or equivalently trace $3$ and determinant $2$.

Conversely, if a real $2\times2$ matrix $B$ has trace $3$ and determinant $2$, its characteristic polynomial is $(t-1)(t-2)$. Choose nonzero eigenvectors $v_1,v_2$ with eigenvalues $1,2$. They are linearly independent: if $v_2=cv_1$, then $2cv_1=Bv_2=cBv_1=cv_1$, forcing $c=0$, a contradiction. The invertible matrix $P=[\,v_1\ v_2\,]$ therefore satisfies $BP=PA$, so $B=PAP^{-1}$. Hence
\[
 \boxed{\operatorname{Orb}(A)=\{B\in M_2(\R):\operatorname{tr}(B)=3,\ \det(B)=2\}.}
\]
All these matrices are invertible because their determinant is $2$.
\end{solution}
\newpage
\begin{questionblock}{5. A normal subgroup of order $5$}
Let $N\trianglelefteq G$, where $|N|=5$ and $|G|$ is odd. Prove that $N\subseteq Z(G)$.
\end{questionblock}
\begin{solution}
Since $|N|=5$, write $N=\langle a\rangle$. Normality makes conjugation by each $g\in G$ an automorphism of $N$, giving a homomorphism
\[
 \theta:G\longrightarrow\Aut(N),\qquad\theta(g)(n)=gng^{-1}.
\]
An automorphism is determined by the image of $a$, which may be any of $a,a^2,a^3,a^4$; each choice does give an automorphism. Thus $|\Aut(N)|=4$. By the first isomorphism theorem and Lagrange's theorem, $|\operatorname{im}\theta|$ divides both $|G|$ and $4$. Since $|G|$ is odd, the image has order $1$. Therefore $gng^{-1}=n$ for all $g\in G$ and $n\in N$, proving $\boxed{N\subseteq Z(G)}$.
\end{solution}
\end{document}
